\chapter{数字经济总论：内涵、特征与测度}

数字经济是全书分析和讲解的对象。本章的任务是把这个对象界定清楚：先给出数字经济的定义与统计口径，再归纳其区别于传统经济的\important{两大成本特征与三大定律}，随后辨析它同信息经济、互联网经济、平台经济、共享经济等相近概念的关系，最后梳理其发展历程与测度方法。全章讨论为后续各章提供对象边界：网络外部性、平台竞争、数据要素、数字金融等主题，都是在数字经济这个界定好的范围内展开的。

\begin{examguide}
\textbf{本章考情}：数字经济的内涵、特征与测度是数字经济专硕各校专业课的必考内容，通常以名词解释、简答题形式出现（\important{专硕·核心}）；数字经济规模测算方法亦见于学硕热点论述题（\important{学硕·热点}）。

\textbf{核心考点}：（1）数字经济定义的演进（Tapscott、OECD、G20 口径）与要素拆解；（2）数字产业化与产业数字化的边界判定；（3）概念谱系六概念的演进关系；（4）三大定律（摩尔、梅特卡夫、吉尔德）的数学表达与经济含义；（5）数字经济测度的两种方法及其局限。

\textbf{本章主线}：什么是数字经济（1.1）→ 它有何特征（1.2）→ 它与相近概念如何区分（1.3）→ 它如何发展而来（1.4）→ 它如何被测量（1.5）。
\end{examguide}

\section{数字经济的内涵与界定}

\subsection{定义的演进}

\textbf{数字经济（Digital Economy）}的概念最早由 Tapscott 于 1996 年在《数字经济：网络智能时代的希望与危险》中提出，用以概括计算机网络所催生的新经济形态。此后国际组织与各国政府相继给出界定，其中两条值得对照记忆。经济合作与发展组织（OECD）在 2020 年发布的数字经济测量框架路线图报告中给出宽口径：数字经济涵盖一切依赖数字投入、或因数字投入而显著增强的经济活动，数字投入包括数字技术、数字基础设施、数字服务与数据。国内考试的标准口径则是 2016 年 G20 杭州峰会通过的《二十国集团数字经济发展与合作倡议》中的定义。

\begin{definition}[数字经济（Digital Economy，G20 杭州峰会，2016）]
数字经济是指以使用数字化的知识和信息作为\important{关键生产要素}、以现代信息网络作为\important{重要载体}、以信息通信技术的有效使用作为\important{效率提升和经济结构优化的重要推动力}的一系列经济活动。
\end{definition}

这一定义以三个并列的"以"字结构，分别规定了数字经济的生产要素、载体与动力。\mn{G20 定义是名词解释的标准答案口径，须逐字记忆；作答时三个"以"缺一不可。}三个短语各有其学理含义。

\textbf{生产要素维度。}"数字化的知识和信息"即数据。传统生产函数以土地、劳动、资本为要素，数字经济将数据确立为关键生产要素，这是生产函数层面的根本变化：数据的非竞争性使规模报酬递增成为可能（详见第 7 章）。此处的"关键"二字并非修辞：它表明数据不是补充性投入，而是决定性的要素。

\textbf{定义的三层结构可以用生产函数来拆解。}将 G20 定义映射到生产函数形式：
\[
Y = A \cdot F\bigl(K, L, D\bigr),
\]
其中数据 $D$ 作为新要素进入要素集合，网络作为载体影响要素的组织方式，信息通信技术的有效使用体现为全要素生产率 $A$。这一映射说明数字经济同时在三个维度上改变生产函数：要素集合（$D$）、要素组织（网络载体）、技术效率（$A$）。任何只强调其一的定义都是不完整的，这也解释了 G20 定义为何以三个"以"字并列展开。

\textbf{载体维度。}"现代信息网络"指以互联网为代表的网络基础设施。要素的流动需要载体：数据只有依托网络才能低成本地复制、传输与汇聚。这一限定同时划出了边界：未依托网络而仅使用单机数字技术的经济活动，严格而言不在定义的完整含义之内，尽管实践中统计口径有所放宽。

\textbf{动力维度。}"信息通信技术的有效使用"被界定为效率提升与结构优化的推动力。限定词"有效使用"蕴含经济学含义：技术本身不创造增长，技术的有效配置与使用才构成生产率提升的来源。第 8 章讨论索洛悖论时还会回到这一点：如果技术普及而未"有效使用"，生产率效应便可能无从显现。

\subsection{统计口径：数字产业化与产业数字化}

在官方统计层面，国家统计局于 2021 年发布《数字经济及其核心产业统计分类（2021）》，将数字经济范围划分为五个大类：数字产品制造业、数字产品服务业、数字技术应用业、数字要素驱动业、数字化效率提升业。\mn{统计分类五大类须记忆：前四类=数字产业化，第五类=产业数字化。2021年文号，考数字年份题常见。}其中前四类构成\textbf{数字产业化}，第五类构成\textbf{产业数字化}。

\begin{definition}[数字产业化与产业数字化]
\textbf{数字产业化（digital industrialization）}指数字技术的产业化部门，即以数字产品与服务为直接产出的行业（统计分类前四类）；\textbf{产业数字化（industrial digitalization）}指数字技术对传统产业的渗透部分，即传统行业使用数字技术所增加的产出（第五类）。
\end{definition}

将数字产业化限定为前四类而非全部五类，依据在于产出的属性。数字产业化的行业具有一个共同特征：\important{其产出本身就是数字产品或数字服务}：芯片、软件、通信设备、互联网信息服务。这些行业离开数字技术则不复存在，故称为"原生"数字部门。\mn{辨析技巧：这个产品离开数字技术还存在吗？存在→产业数字化，不存在→数字产业化。}产业数字化则不同：制造业、农业、金融业的产出在数字技术出现之前已然存在，数字技术的作用是改造生产过程、提升效率。因此产业数字化的测度不能直接取传统行业的全部增加值，而需剥离出其中"由数字技术贡献的部分"，这正是本章末节测度难点的由来。

\begin{misconception}
\textbf{误区}："数字化效率提升业"对应的产业数字化，不能将传统行业全部增加值计入数字经济规模。若将使用数字技术的企业全部产出计入，则数字经济规模将与整个国民经济规模趋同，测度失去意义。正确的做法是仅计入数字技术贡献的增量部分（实践中以数字化投入比例或行业数字化渗透率进行剥离）。
\end{misconception}

统计口径的分歧在中美数字经济的规模对比中表现得最典型。2020 年中国数字经济规模约 39.2 万亿元、占 GDP 约 38.6\%（中国信息通信研究院《中国数字经济发展白皮书》，2021），同年美国数字经济规模约 2.14 万亿美元、占 GDP 约 10.2\%（美国经济分析局，2021）。两组数据是否说明中国数字经济发展程度远高于美国？答案是否定的，差异主要来自口径：美国经济分析局（BEA）的测算只覆盖数字经济的核心产业，大致对应我们的数字产业化；信通院口径则在核心产业之外加计产业数字化的渗透贡献。\mn{中美占比辨析：中国 38.6\% 是"宽口径"（含产业数字化），美国 10.2\% 是"窄口径"（仅核心产业），两者不可直接比较。比较数字经济规模时必须先对齐口径。}产业数字化的体量远大于核心产业，宽口径自然远高于窄口径。\textbf{口径差异正是本章 1.5 节测度难题的总根源}：数字经济的边界划在哪里，规模就量到哪里，不同机构对"数字技术贡献的剥离比例"认定不同，同一国家的数字经济规模可以相差数倍。

\subsection{"四化"框架}

统计分类回答的是"数字经济规模怎么量"的问题，政策推进还需要回答"数字经济怎么发展"的问题。为此，中国信通院在统计口径的基础上给出"四化"框架，把数字经济的政策任务概括为四个相互关联的部分。

\begin{definition}[数字经济"四化"框架（信通院）]
中国信通院将数字经济划分为四个部分：数字产业化、产业数字化、数字化治理、数据价值化。前两化对应统计口径的核心部分，后两化是政策层面的延伸。
\end{definition}

\begin{itemize}
\item \textbf{数字产业化}：信息通信产业本身（前述五大类中前四类）；
\item \textbf{产业数字化}：传统产业应用数字技术带来的产出增加与效率提升（第五类）；
\item \textbf{数字化治理}：政府运用数字技术进行治理（数字政府），以及对数字经济的治理（平台监管）；前者如政务服务"一网通办"，后者如平台反垄断与数据安全监管；
\item \textbf{数据价值化}：数据资源化、资产化与资本化的过程，即数据从原始记录转化为可定价、可交易、可入表的资产；《关于构建数据基础制度更好发挥数据要素作用的意见》（"数据二十条"，2022 年）是这一过程的制度基础（详见第 7 章）。
\end{itemize}

"四化"之间的关系是：数字产业化提供技术与产品供给，产业数字化创造需求与应用场景，数字化治理提供制度环境，数据价值化贯穿其中，使前两化产生的数据沉淀为可再度投入生产的新型要素。

\section{数字经济的特征与三大定律}

1.1 节回答了数字经济是什么。本节回答数字经济与传统经济的根本区别：数字产品的成本结构、数字经济的报酬性质，以及三条刻画数字技术演进的定律。理解这些特征是后续平台、数据、增长等各章分析的共同前提。

\subsection{零边际成本}

数字产品与传统产品在生产成本结构上存在根本差异。数字产品（软件、数据、数字内容）的生产以高额固定成本（研发、内容制作）为代价，而其复制与分发的边际成本趋近于零。

\begin{derivation}{零边际成本与竞争定价失效}{}
\textbf{零边际成本下的定价问题。}设数字产品生产的固定成本为 $F$，边际成本为 $c \to 0$。在完全竞争市场中，厂商是价格接受者，均衡价格趋于边际成本 $P = c \approx 0$，厂商无法回收固定成本 $F$，市场供给将不存在。形式化地，给定市场价格 $P$，厂商利润为
\[
\pi(q) = P \cdot q - F - cq = (P-c)\,q - F,
\]
当 $c \to 0$ 且竞争使价格趋于边际成本（$P \to c \approx 0$）时，$\pi \to -F < 0$：任何正产量均导致亏损，厂商的最优选择是不供给。因此\important{零边际成本产品与完全竞争定价逻辑不相容}。
\end{derivation}

这一结论的经济含义是：完全竞争定价（$P=MC$）的前提是边际成本为正、厂商通过产量竞争实现配置效率；当 $MC\approx 0$ 时，这一机制失效。零边际成本行业必须依赖新的定价机制（版本区分、订阅制、广告补贴、捆绑销售等，第 5 章的主题）回收固定成本，故其天然具有垄断或差异化定价的倾向。这一结论为平台经济的市场结构分析（第 4、6 章）提供了成本侧的基础。

零边际成本的经典案例是维基百科对大英百科全书的替代。\mn{案例：维基百科（开放编辑、零边际成本分发）在 2012 年终结了 244 年历史的大英百科全书纸质版。经济学含义：零边际成本 + 开放协作将内容生产的固定成本也外部化了。}2005 年《自然》杂志发表的专家对比研究显示，维基百科科学条目的准确性与大英百科全书接近，而前者对读者完全免费。维基百科同时将固定成本（条目撰写）以开放协作的方式外部化，边际成本（一次页面访问的服务器开销）近乎为零。当一种产品的生产与分发成本同时趋于零时，基于"稀缺性+排他性"的\important{传统内容商业模式便不再成立}，这正是零边际成本特征在现实中最完整的体现。

\subsection{边际收益递增}

传统工业经济服从边际收益递减：在技术不变、其他要素投入固定的条件下，持续增加单一要素投入，其边际产出终将下降。这一规律是传统生产函数凹性的直接体现：若 $Y=F(K,L)$ 满足 $F_{KK}<0$、$F_{LL}<0$，则每种要素的边际产出随投入增加而递减，企业的最优规模由此内生地确定。需要说明的是，严格的经济学教科书称这一规律为\textbf{边际报酬递减}（或规模报酬递减），"边际收益"在微观经济学中原指边际收入；数字经济文献与考试大纲通行的"边际收益递增"说法，实指边际报酬递增。本书与考试口径一致，沿用"边际收益"的表述。

数字经济至少在两个机制上逆转了这一逻辑。其一，\textbf{网络外部性使需求侧报酬递增}：用户规模越大，产品对每个用户的价值越高，愿意加入的用户越多，这是"越用越有价值"的需求侧效应，与供给侧的投入产出关系无关（下节梅特卡夫定律将给出 $V\propto n^2$ 的数理形式）。其二，\textbf{数据要素的非竞争性使供给侧报酬递增}：资本与劳动在使用中被消耗或占用，而数据可被多个生产过程同时使用且不被损耗；随着数据规模与多样性增加，其可挖掘的信息价值反而上升（第 7 章）。\mn{边际收益递增记忆两点：需求侧靠网络外部性，供给侧靠数据非竞争性；后果：凸性假设失效，反垄断分析变难。}二者叠加，数字产品的生产与消费均呈现"越大越强"的正反馈。递增报酬对市场结构的影响，可以用一个简单模型刻画。

\begin{derivation}{规模报酬与竞争均衡的存在性}{}
\textbf{设定。}设生产函数为齐次函数，满足 $f(\lambda K, \lambda L) = \lambda^{\gamma} f(K,L)$，其中 $\gamma$ 刻画规模报酬性质：$\gamma<1$ 为规模报酬递减，$\gamma=1$ 为规模报酬不变，$\gamma>1$ 为规模报酬递增。竞争厂商以市场价格 $P$ 出售产出、以工资 $w$ 与利率 $r$ 购买要素，利润为
\[
\pi(K,L) = P \cdot f(K,L) - wK - rL.
\]

\textbf{基准情形：规模报酬递减（$\gamma<1$）。}固定要素投入比例 $(K,L)$，将利润视为规模倍数 $t$ 的函数，由齐次性得
\[
\pi(tK, tL) = P \cdot t^{\gamma} f(K,L) - t\,(wK + rL).
\]
当 $t>1$ 且 $\gamma<1$ 时 $t^{\gamma}<t$，故
\[
\pi(tK, tL) < t\,[P \cdot f(K,L) - wK - rL] = t\,\pi(K,L),
\]
即放大规模时利润的增长慢于规模本身的增长。进一步，当 $t\to\infty$ 时 $t^{\gamma-1}\to 0$，成本项 $t(wK+rL)$ 的增长快于收入项 $P\,t^{\gamma} f(K,L)$，$\pi(tK,tL)\to-\infty$：无限扩大规模终将亏损，利润函数沿规模方向先升后降，存在有限的利润最大规模。该规模由一阶条件确定：
\[
\gamma P\,t^{\gamma-1}f(K,L) - (wK+rL) = 0,
\]
即沿该投入射线，边际产出恰好等于要素成本总额；二阶条件 $\gamma(\gamma-1)P\,t^{\gamma-2}f(K,L) < 0$ 保证其为唯一的最大值点。因此\important{规模报酬递减（正常情形）下，利润最大化有有限的内点解，竞争均衡存在，企业最优规模内生确定}。作为参照，$\gamma=1$（规模报酬不变）时 $\pi(tK,tL) = t\,\pi(K,L)$，利润随规模等比例放大，零利润条件虽可维持，但厂商规模不确定。

\textbf{数字经济情形：规模报酬递增（$\gamma>1$）。}若在 $(K,L)$ 处利润 $\pi(K,L) \ge 0$，则对任意放大倍数 $t>1$，由齐次性有
\[
\pi(tK, tL) = P \cdot t^{\gamma} f(K,L) - t\,(wK + rL) > t\,\pi(K,L) \ge 0.
\]
即放大生产规模总能提高利润：若 $\pi(K,L) > 0$，利润随规模无限放大，不存在有限的最优规模；若 $\pi(K,L) = 0$，放大后 $\pi(tK,tL) > 0$，零利润的竞争均衡条件无法维持。因此\important{规模报酬递增破坏竞争均衡的存在性条件}：传统竞争均衡赖以成立的凸性假设不再满足。
\end{derivation}

两种情形的对照正是本节的考点：\important{规模报酬递减（正常情形）下生产集凸、均衡存在、企业最优规模内生确定；规模报酬递增时凸性失效，均衡可能不存在、不唯一或不稳定，市场可能收敛于垄断而非竞争}。这正是数字经济反垄断分析中市场界定与市场势力评估难题的理论根源（第 6 章）。

\subsection{三大定律}

三条经验定律分别刻画数字技术扩张的三个侧面：摩尔定律对应计算成本，梅特卡夫定律对应网络价值，吉尔德定律对应带宽成本。三者共同解释了数字技术为何能以不断下降的成本支撑不断扩张的规模。

\begin{definition}[摩尔定律（Moore's Law）]
集成电路上可容纳的晶体管数目约每 18--24 个月翻一番。
\end{definition}

该定律由英特尔创始人戈登·摩尔于 1965 年基于经验观察提出，\mn{摩尔定律是经验观察而非物理定律，其持续依赖于研发投入的指数增长；近年已出现放缓迹象，考试辨析题常考"定律"与"规律"之别。}并预测这一趋势将持续。

\begin{derivation}{摩尔定律的数学表达}{}
设 $N(t)$ 为 $t$ 时刻晶体管数目，翻番周期为 $T$（取 18 个月），则
\[
N(t) = N_0 \cdot 2^{\,t/T}.
\]
两边取自然对数：
\[
\ln N(t) = \ln N_0 + \frac{t}{T}\ln 2,
\]
即 $\ln N$ 对 $t$ 呈线性增长，斜率为 $\ln 2 / T$。这一线性化形式给出了验证摩尔定律的计量方法：若数据在半对数坐标系中呈直线，则指数增长假设成立。进一步，可以导出摩尔定律的价格含义。历史经验表明，单颗芯片的生产成本大体保持稳定：设其为常数 $C_0$，则单位计算能力（单个晶体管）的价格为
\[
p(t) = \frac{C_0}{N(t)} = \frac{C_0}{N_0} \cdot 2^{-t/T} = p_0 \cdot 2^{-t/T},
\]
其中 $p_0 \equiv C_0 / N_0$ 为基期单位计算能力价格：芯片成本不变而晶体管数翻倍，单位计算能力的价格便按同一周期减半。因此\important{单位计算能力价格以指数速度衰减}。
\end{derivation}

计算能力价格指数衰减意味着，同等预算所能购买的计算能力以指数速度增长。这是\important{摩尔定律真正的经济含义}：它使数字技术向各行业的渗透在成本上持续可行，构成数字经济扩张的技术前提；同时也解释了为何数字技术会不断"侵蚀"那些曾被认为不会数字化的行业：当计算成本每一年半减半时，几乎所有依赖信息处理的环节终将被自动化覆盖（第 9 章就业冲击的技术根源）。

摩尔定律刻画的是供给侧的算力成本，梅特卡夫定律转向需求侧：网络的价值不取决于网络本身，而取决于使用网络的人以及人与人之间的连接。

\begin{definition}[梅特卡夫定律（Metcalfe's Law）]
网络的价值与用户数量的平方成正比，即 $V \propto n^2$。
\end{definition}

该定律由 3Com 公司创始人梅特卡夫提出，其推导如下。

\begin{derivation}{梅特卡夫定律的推导与修正}{}
考虑一个具有 $n$ 个用户的网络。若每个用户与其他任一用户之间构成一条潜在连接，则连接总数为组合数
\[
C_n^2 = \frac{n(n-1)}{2} = \frac{n^2 - n}{2}.
\]
当 $n$ 较大时，$n^2 - n \approx n^2$，故网络价值 $V$ 满足
\[
V \propto n^2.
\]
这一推导的假设是：\important{每条连接的潜在价值相同}。Briscoe 等学者质疑该假设，指出连接价值服从齐普夫分布——第 $k$ 有价值的连接，其价值约为 $1/k$，少数连接贡献绝大部分价值。此时单个用户的 $n-1$ 条连接的价值之和近似为调和级数 $\sum_{k=1}^{n-1} \frac{1}{k} \approx \ln n$，$n$ 个用户的总价值因此为
\[
V \approx n \cdot \ln n,
\]
即 $V \propto n \ln n$。
\end{derivation}

这一结果的直觉在于网络价值的加速增长：按组合数口径，第 $n+1$ 个用户加入带来的边际价值约为 $n$，且随用户规模增大而递增，因为他同时增加了所有存量用户的可连接对象。这种边际价值递增是网络型产品与传统产品的本质区别，也是平台补贴获客策略（第 4、5 章）的价值论依据：吸引一个用户的收益并非来自该用户本身，而是来自他提升的所有存量用户的价值。两种口径的差异同样具有实质含义：按 $n^2$ 口径，网络价值的增速快于用户规模增速本身；按 $n\ln n$ 口径，网络价值增速仅略高于线性，补贴的合理性需要更谨慎的论证。\mn{梅特卡夫定律标准表述为 $V\propto n^2$，$n\ln n$ 为修正口径。}掌握两种口径的差异之后，考试作答仍以标准表述 $V\propto n^2$ 为准。

\begin{misconception}
\textbf{误区}：梅特卡夫定律刻画的是网络的\important{潜在}价值（连接的可能性），而非企业实际实现的收入。网络价值 $n^2$ 与平台利润之间隔着\important{商业模式（变现方式）}，二者不可等同；将梅特卡夫定律直接用作平台估值的依据，是分析中的常见错误：2000 年前后互联网泡沫的估值狂热，很大程度上源于对"潜在价值"与"实现收入"的混淆（见 1.4.2 节案例）。
\end{misconception}

\begin{definition}[吉尔德定律（Gilder's Law）]
主干网带宽的增长速度约每 6 个月翻一番，约为摩尔定律（18 个月）速度的三倍。
\end{definition}

该定律由吉尔德于 1996 年提出。与摩尔定律同理，带宽增长可写为指数形式 $B(t) = B_0 \cdot 2^{\,t/T_B}$，其中翻番周期 $T_B = 6$ 个月。

\begin{derivation}{吉尔德定律与摩尔定律的速度比较}{}
设带宽翻番周期为 $T_B$，算力翻番周期为 $T_M$。对指数增长函数 $N(t)=N_0 2^{t/T}$ 求对数增长率：
\[
\frac{d \ln N(t)}{dt} = \frac{\ln 2}{T},
\]
即翻番周期越短，对数增长率越高，两者成反比。取 $T_B = 6$ 个月、$T_M = 18$ 个月：
\[
\frac{g_B}{g_M} = \frac{\ln 2 / T_B}{\ln 2 / T_M} = \frac{T_M}{T_B} = 3,
\]
即带宽的对数增长率约为算力的三倍。
\end{derivation}

吉尔德据此推论"带宽终将免费"，其经济含义在于：当带宽的边际成本趋于零时，依托带宽传输的商业模式（视频、云计算）的扩张约束被解除，为"免费+增值"商业模式（第 5 章）提供了基础设施条件。带宽与算力的双指数下降共同解释了数字经济扩张的基础设施逻辑：传输成本下降决定用户数量"连得进"，计算成本下降决定"算得起"。

\begin{figure}[htbp]
\centering
\begin{tikzpicture}[>=Stealth, scale=0.9]
\begin{axis}[
    width=10.5cm, height=6cm,
    xlabel={用户数 $n$}, ylabel={网络价值 $V$},
    xmin=0, xmax=20, ymin=0, ymax=220,
    legend pos=north west,
    grid=major,
    axis lines=left,
]
\addplot[color=main, very thick, domain=0:20, samples=100] {x^2/2};
\addlegendentry{$V \propto n^2$（梅特卡夫）}
\addplot[color=winered, thick, domain=0:20, samples=100] {x*ln(x+1)*2.085};
\addlegendentry{$V \propto n\ln n$（修正）}
\addplot[color=third, thick, domain=0:20, samples=100] {5*x};
\addlegendentry{$V \propto n$（线性）}
\end{axis}
\end{tikzpicture}
\caption{梅特卡夫定律及其修正形式的比较}
\label{fig:metcalfe}
\end{figure}

图~\ref{fig:metcalfe} 对比了三种口径下网络价值随用户规模的增长路径。为便于比较，三条曲线在 $n=10$ 处取值相同；在此基准点右侧，$n^2$ 口径的价值增速最快，与另两条曲线的差距随用户规模迅速拉大，这是"临界规模"概念的图形基础（第 3 章）：在网络价值加速增长与用户采用相互强化的正反馈下，平台一旦越过临界规模便难以被后来者超越。

\subsection{特征总结}

表~\ref{tab:features} 将前述差异归纳为六个维度。

\begin{table}[H]
\centering
\caption{数字经济与传统经济的特征对照}
\label{tab:features}
\begin{tabular}{p{2.6cm}p{5.2cm}p{4.4cm}}
\hline
\textbf{维度} & \textbf{传统工业经济} & \textbf{数字经济} \\
\hline
核心生产要素 & 土地、劳动、资本 & 数据（+算法、算力） \\
成本结构 & 边际成本先降后升、最低点为正 & 复制边际成本趋于零 \\
规模报酬 & 递减为主 & 递增（网络外部性、数据） \\
竞争形态 & 原子式竞争或供给方规模垄断 & 需求方规模经济、赢家通吃 \\
定价逻辑 & 边际成本定价 & 免费、补贴、歧视性定价 \\
市场边界 & 地理与产品边界清晰 & 跨界竞争、平台包抄 \\
\hline
\end{tabular}
\end{table}

表~\ref{tab:features} 中"数据（+算法、算力）"的写法对应政策话语中的数字经济\textbf{三要素}：数据是原材料，算法是加工方法，算力是加工能力，三者缺一不可。\mn{三要素记忆：数据是"料"、算法是"方"、算力是"力"，对应生产中的原料、工艺与产能三个环节。}没有算法与算力，数据只是未经处理的原始记录；没有数据，算法与算力便无事可做。

\section{相关概念辨析}

数字经济的相关概念繁多，考试中辨析题的命题点集中于两组关系：一是谱系关系（信息经济、互联网经济、网络经济、平台经济、共享经济与数字经济的关系），二是口径关系（数字产业化与产业数字化的关系，见 1.1.2 节）。

\subsection{概念的演进谱系}

图~\ref{fig:lineage} 给出了数字经济相关概念的演进路径。概念谱系的逻辑是：技术形态的递进（计算机、互联网、平台、智能）决定了经济形态的递进，后一个概念以前一个概念为基础并扩展其范围。

\begin{figure}[htbp]
\centering
\begin{tikzpicture}[>=Stealth,
    box/.style={rectangle, rounded corners=2mm, draw=second, thick, fill=second!8, align=center, font=\small, minimum height=9mm, minimum width=22mm},
    arr/.style={->, thick, main}]
\node[box] (A) at (0,0) {信息经济\\(1960s--)};
\node[box] (B) at (45mm,0) {互联网经济\\(1990s--)};
\node[box] (C) at (90mm,0) {网络经济\\(1990s--)};
\node[box] (D) at (90mm,-16mm) {平台经济\\(2000s--)};
\node[box] (E) at (45mm,-16mm) {共享经济\\(2010s--)};
\node[box] (F) at (0,-16mm) {数字经济\\(1996--)};
\node[anchor=north, font=\small\color{structurecolor}] at (A.south) {计算机普及};
\node[anchor=north, font=\small\color{structurecolor}] at (B.south) {网络连接};
\node[anchor=north, font=\small\color{structurecolor}] at (C.south) {网络效应};
\node[anchor=north, font=\small\color{structurecolor}] at (D.south) {双边市场};
\node[anchor=north, font=\small\color{structurecolor}] at (E.south) {使用权交易};
\node[anchor=north, font=\small\color{structurecolor}] at (F.south) {数据要素};
\draw[arr] (A) -- (B);
\draw[arr] (B) -- (C);
\draw[arr] (C) -- (D);
\draw[arr] (D) -- (E);
\draw[arr] (E) -- (F);
\end{tikzpicture}
\caption{数字经济相关概念的演进谱系}
\label{fig:lineage}
\end{figure}

各概念之间的关系可以概括为四点。\textbf{第一，信息经济是源头}：马克卢普与波拉特在 20 世纪 60--70 年代提出信息部门分类，是最早将"信息的产业化"纳入经济分析框架的尝试。\textbf{第二，互联网经济与网络经济是中间形态}：二者都以网络连接为基础，互联网经济侧重网络化的商业活动（电子商务、门户网站），网络经济侧重网络效应对资源配置方式的改变，两者大体同时出现，常被混用。\textbf{第三，平台经济与共享经济是典型实现}：平台经济以双边市场为组织形态，是数字经济的组织层；共享经济以闲置资源使用权交易为特征，是平台经济在生活服务领域的细分形态。\textbf{第四，数字经济是总括}：数字经济以数据要素与 ICT 技术为统一内核，将前述各形态纳入同一分析框架。\mn{谱系题口诀：信息是源头、网络是连接、平台是组织、共享是细分、数字是总括。辨析时先判断概念对应哪一阶段的技术形态。}概括起来：\important{信息经济是源头，互联网经济与网络经济是中间形态，平台经济与共享经济是典型实现，数字经济是总括形态。}

此外，\textbf{数字化、数字化转型与智能化}构成企业层面的递进关系：数字化指业务信息的数字记录（信息化），数字化转型指以数字技术重构业务流程与商业模式，智能化指以人工智能替代或增强决策环节。三者的递进对应着数字技术从"记录工具"到"生产方式"再到"决策主体"的角色跃迁。

\section{数字经济发展历程与格局}

\subsection{三个阶段}

数字经济的发展可划分为三个阶段，其划分依据是主导技术的形态。这一划分与 1.3 节的概念谱系互相印证：每一阶段的主导技术，恰好对应一个或几个新经济概念的兴起。三个阶段并非相互取代，而是逐层叠加，上一阶段积累的要素与能力构成下一阶段的起点。

\textbf{信息化阶段（1940s--1990s）。}以计算机与通信技术的发明与普及为主线。此阶段数字技术的角色是"记录与计算工具"：信息的存储、计算与传输成本大幅下降，为后续阶段的要素积累奠定技术基础。

\textbf{网络化阶段（1990s--2010s）。}以互联网的商业化为起点。此阶段数字技术的角色是"连接与组织工具"：连接的成本趋近于零使市场边界急剧扩张，电子商务、搜索引擎、社交网络相继兴起，平台作为新组织形态登上经济舞台。

\textbf{智能化阶段（2010s--）。}以大数据、云计算与人工智能的成熟为标志。此阶段数字技术的角色是"要素与决策工具"：数据从生产过程的副产品转变为核心生产要素，算法从辅助工具转变为配置资源的主体，数字经济对增长、就业与分配的影响全面显现。

\subsection{案例：互联网泡沫与网络价值的第一次集体误判}

网络化阶段留给经济学最重要的历史教训，是 1995--2000 年的互联网泡沫。\mn{案例：纳斯达克指数 1995 年约 1000 点，2000 年 3 月峰值 5048 点，此后两年下跌约 78\%。泡沫期间大量"概念公司"并无收入，估值依据是"眼球数"与"点击量"。}泡沫期间，投资者将梅特卡夫定律误读为估值公式：既然网络价值与用户数的平方成正比，那么任何用户数快速增长的公司似乎都值高价。这一误读混淆了两个概念：网络潜在价值的增速（$n^2$）与公司实现收入的能力（取决于变现方式），其间的差别可以用价值函数精确刻画。

\begin{derivation}{网络价值与变现价值}{}
设网络用户数为 $n$，网络的\textbf{潜在价值}（连接价值）为 $V(n) = \kappa n^2$（梅特卡夫口径）。公司的\textbf{实现收入}取决于变现率，即每个用户产生的平均收入 $\mathrm{ARPU}$，它由商业模式的转化效率 $\rho$ 与用户支付能力共同决定：
\[
R(n) = \rho \cdot \mathrm{ARPU} \cdot n.
\]
泡沫期的估值错误在于将 $V(n)$ 直接当作 $R(n)$ 使用。定义"价值实现率"为
\[
\lambda(n) = \frac{R(n)}{V(n)} = \frac{\rho \cdot \mathrm{ARPU}}{\kappa n}.
\]
$\lambda$ 随用户规模\textbf{递减}：潜在价值以 $n^2$ 增长，而变现收入以 $n$ 增长，用户规模越大，价值实现的缺口越大。
\end{derivation}

泡沫公司的问题恰在于此：以 $V(n)$ 的高增长故事融资，却以 $R(n)$ 的低斜率现实交付。变现收入以 $n$ 而非 $n^2$ 增长，原因在于变现受限于用户的时间与支付能力。泡沫破裂后幸存的公司（亚马逊、谷歌）的共同特征，是找到了使 $\rho$ 持续上升的变现机制（亚马逊的电商闭环、谷歌的搜索广告）。\important{网络效应是价值创造的必要条件而非充分条件，变现效率 $\rho$ 的持续提升才是价值实现的最终约束}。

\subsection{全球格局：中美两强与主要经济体战略}

从国际比较看，全球数字经济呈现"中美两强"格局：美国在底层技术与平台企业的全球影响力上领先，中国在市场规模、产业应用与用户规模上领先，两国合计构成全球数字经济的主要份额。\mn{全球格局答题口径：美国领先底层技术与平台全球化，中国领先市场规模与产业应用；论述题引用信通院"全球第二"的定位即可。}其他主要经济体也各有布局：欧盟以《数字十年》规划提出 2030 年数字目标，德国推动"工业 4.0"，日本提出"社会 5.0"，其共同点是都试图在制造业或社会治理领域确立数字技术的先发优势。

\subsection{中国数字经济发展概况与政策脉络}

按中国信通院测算，2023 年中国数字经济规模达 53.9 万亿元，占 GDP 比重为 42.8\%。\mn{数据口径：53.9 万亿元、42.8\%（信通院，2023 年）。此类数据用于论述题背景。}这一规模的国际比较地位（总量全球第二、增速快于 GDP 增速）是论述题中"数字经济成为增长新引擎"命题的实证支撑。

政策脉络方面，2023 年中共中央、国务院印发的《数字中国建设整体布局规划》提出"2522"整体框架，是当前政策表述的标准口径：

\begin{itemize}
\item \textbf{两大基础}：数字基础设施（算力、网络）、数据资源体系（数据要素市场）；
\item \textbf{五位一体}：数字技术与经济、政治、文化、社会、生态文明建设五位一体深度融合；
\item \textbf{两大能力}：数字技术创新体系、数字安全屏障；
\item \textbf{两个环境}：国内数字治理生态、国际数字合作环境。
\end{itemize}

"2522"框架的提出有其现实针对性：此前数字经济政策分散于产业政策、网络治理、数据立法等条线，缺乏跨部门的统筹安排；平台经济经历高速扩张与规范整顿后，也需要一个"发展与规范并重"的总体设计。\mn{"2522"框架为 0258 简答题标准素材：2 基础 + 5 融合 + 2 能力 + 2 环境。记忆线索：基础、融合、能力、环境，由内而外。}这一布局的经济含义在于，数字基础设施与数据资源体系直接扩大数字经济的供给能力（算力与数据是新型要素投入），五位一体融合决定数字技术向各领域的渗透广度，两个环境则提供制度与国际空间的保障；"基础、融合、能力、环境"由内而外的结构，恰好对应 1.5 节测度框架中"要素、应用、环境"的递进逻辑。

\section{数字经济测度}

数字经济规模不可直接观察，必须经由统计方法测度。现行主流方法分为增加值测度法与指数测度法两大体系：前者回答"数字经济规模有多大"（绝对量），后者回答"数字经济发展有多快"（相对变化）。

\subsection{增加值测度法}

增加值测度法的理论根基是国民经济核算恒等式。数字经济规模不可直接观察，只能借助国民经济核算的既有框架，把"数字部门"从国民经济中圈定并剥离出来。下面的推导说明这一圈定与剥离如何落实到公式层面。

\begin{derivation}{生产法增加值恒等式}{}
对任一生产单位，其\textbf{增加值}定义为总产出与中间投入之差：
\[
VA = P \cdot Q - \sum_i p_i m_i,
\]
其中 $Q$ 为产出量，$P$ 为产出价格，$m_i$ 为第 $i$ 种中间投入量，$p_i$ 为其价格。国民经济总量等于各部门增加值之和：
\[
GDP = \sum_j VA_j.
\]
数字经济增加值测度的具体步骤为：\textbf{第一步}，依《数字经济及其核心产业统计分类》从国民经济行业分类中筛选数字产业化行业；\textbf{第二步}，按生产法加总这些行业的增加值，得到数字产业化规模；\textbf{第三步}，对产业数字化部分，按数字化投入比例（$s_d$）剥离传统行业增加值中的数字贡献，得到产业数字化规模；\textbf{第四步}，两者相加即数字经济总规模：
\[
\text{数字经济规模} = \underbrace{\sum_{j \in \text{核心产业}} VA_j}_{\text{数字产业化}} + \underbrace{\sum_{k \in \text{传统产业}} s_d^k \cdot VA_k}_{\text{产业数字化}}.
\]
\end{derivation}

增加值测度以生产法为基础，但国民经济核算恒等式还给出另外两条等价路径，它们对数字经济的间接测算与交叉验证不可或缺。收入法从分配角度加总生产所创造的要素收入，支出法从需求角度加总对最终产品的购买，三条路径在核算上严格相等，构成测度方法相互校验的理论基石。

\begin{derivation}{收入法与支出法恒等式}{}
\textbf{收入法}从增加值分配的角度核算：一国生产创造的价值全部以要素报酬的形式流向居民、政府与企业，
\[
GDP = \sum_j \big(\text{劳动者报酬}_j + \text{生产税净额}_j + \text{固定资产折旧}_j + \text{营业盈余}_j\big),
\]
其中生产税净额 = 生产税 $-$ 补贴，营业盈余为扣除前三项后归企业支配的剩余。\textbf{支出法}从最终使用的角度核算：所有最终产品的去向只有消费、投资、政府购买与净出口四类，
\[
GDP = C + I + G + (X - M),
\]
其中 $C$ 为居民消费、$I$ 为资本形成总额、$G$ 为政府购买、$X-M$ 为净出口。三种方法核算的是同一对象，恒等式
\[
\sum_j VA_j = \sum_j \big(\text{劳动者报酬}_j + \text{生产税净额}_j + \text{折旧}_j + \text{营业盈余}_j\big)
\]
与
\[
\sum_j VA_j = C + I + G + (X - M)
\]
同时成立。这一恒等式对数字经济的意义在于：生产法要求行业增加值数据，支出法只要求最终产品的支出分类数据，收入法只要求要素收入数据。当生产侧数据不可得（如部分数字产品未进入行业统计）时，可从支出侧或收入侧反向推算数字经济的规模，并以三法互验。
\end{derivation}

增加值度量的是一国（或一行业）\important{新创造}的价值：用总产出减去中间投入，排除了从其他生产单位购入的投入品价值，从而避免重复计算。数字经济测度无非是将这一国民经济核算的基本原则应用于一个新界定的部门：先圈定"原生数字部门"（前四类）直接加总，再对"被数字技术改造的传统部门"按比例剥离贡献。方法本身并不新颖，新颖（也困难）的是圈定与剥离这两个环节。\mn{增加值法四步：筛行业、加总前四类、剥离第五类、求和。推导题可能要求写出上述测算公式并说明 $s_d$ 的含义。}

由此可得增加值测度的两条操作路径。上述四步属于\textbf{直接法}，即生产法：从生产侧直接核算数字部门的增加值。与之相对的是\textbf{间接法}，即需求法：从支出侧入手，通过数字相关消费、投资与出口等支出项估算数字经济规模；当生产侧基础数据不足时，间接法可作为补充或交叉验证。收入法在数字经济测度中主要用于两处：一是三法互验，识别并修正生产法中的重复计算或遗漏；二是作为增长核算法（第 8 章）的理论基础，增加值按要素贡献分解正是索洛框架的起点。三种路径的差异不在结论而在数据可得性：直接法依赖行业增加值数据，间接法依赖支出分类数据，收入法依赖要素收入与折旧数据。

增加值测度法面临三类公认的测度难题。\textbf{其一，免费数字产品的价值核算}：搜索引擎、社交平台对消费者免费提供，其价值通过广告收入间接实现，但消费者剩余部分不计入 GDP，\important{导致数字经济规模被系统性低估}（这一低估的方向性论证见第 8 章对索洛悖论的分析）。\textbf{其二，质量调整问题}：数字产品的性能快速提升伴随价格下降，若价格指数未充分进行质量调整，实际产出将被低估。\mn{测度难点三件套：免费产品、质量调整、剥离比例与重复计算。论述题"数字经济测度的挑战"按此展开。}\textbf{其三，重复计算与剥离比例问题}：数字部门与传统部门互为投入，直接加总可能重复计算；而产业数字化的剥离比例 \important{$s_d$} 缺乏公认标准，不同机构对数字化投入份额的认定差异显著，导致同一国家的数字经济规模在不同口径下差异可观。分歧的规模有直接证据：2016 年中国数字经济占 GDP 比重，不同机构的测算结果从约 5.6\%（布鲁盖尔研究所）到约 30.6\%（腾讯研究院）不等；美国的同一指标，从约 5.4\% 到约 58.3\% 的区间同样悬殊。\mn{测度分歧的实证锚点：同一国家同一年的数字经济占比，不同机构可以相差数倍（中国约 5.6\%--30.6\%，美国约 5.4\%--58.3\%，2016 年）。论述题"为什么数字经济测不出唯一答案"直接用这组数据开头。}差异的来源正是剥离比例、免费产品处理与质量调整三件难点的不同技术选择，这意味着\important{任何数字经济规模数据都必须连同口径一起引用}。

\subsection{指数测度法}

指数测度法将多维度指标合成为单一指数，用于刻画数字经济的发展水平与动态。当问题从"规模有多大"转向"发展有多快、水平孰高孰低"时，单一指标无法回答，需要把多个维度的信息压缩成一个可比数字。合成过程分两步：标准化与加权。

\begin{derivation}{指数的标准化与合成}{}
设待合成的一级指数由 $m$ 个二级指标构成，指标 $j$ 的原始取值为 $x_j$。\textbf{第一步，无量纲化}：由于各指标量纲与量级不同，直接加权无意义，需先标准化。常用极差法：
\[
x_j' = \frac{x_j - x_{\min}}{x_{\max} - x_{\min}} \in [0,1].
\]
\textbf{第二步，加权合成}：设指标权重为 $w_j$（$\sum_j w_j = 1$），则合成指数为
\[
I = \sum_{j=1}^{m} w_j x_j'.
\]
权重的确定是合成指数的核心环节：层次分析法（AHP）以专家两两比较构造判断矩阵，通过特征向量确定权重；熵值法以数据自身的离散程度（信息熵）客观赋权（推导详见第 13 章）。
\end{derivation}

指数法的本质是将"发展水平"这一不可直接观测的概念操作化为可计算的多维加权平均。极差标准化只保留指标的相对位置信息（最高为 1、最低为 0），加权合成则体现各维度的相对重要性。因此指数的可靠性完全取决于两件事：指标体系是否完备（是否漏掉了重要维度）、权重方案是否合理（主观赋权与客观赋权的取舍）。

中国信通院于 2017 年发布的数字经济指数（DEI）是当前引用率最高的数字经济指数之一。DEI 属景气指数，由先行、一致、滞后三类指数构成，其指标体系从宏观经济、基础能力、基础产业、融合应用四个层面共选取 23 个指标（图~\ref{fig:dei} 为示意图，二级节点列示各层面的代表性指标）。此外，各省级统计部门与学术机构亦发布地区性数字经济指数，其构造逻辑与上述两步（标准化+加权）一致。

\important{【专硕·扩展】}在国内指数体系中，除信通院 DEI 外，财新智库等机构发布的中国数字经济指数也被广泛引用；两者侧重不同，DEI 属景气指数，中国数字经济指数更偏重产业发展与融合应用维度。学界对数字经济发展效率的测度常借用\textbf{DEA-Malmquist} 方法：数据包络分析（DEA）以线性规划构造生产前沿面，度量各决策单元相对前沿的效率；Malmquist 指数进一步将效率变化分解为技术效率变化与技术进步，用于刻画数字经济的效率演进。

\begin{figure}[htbp]
\centering
\begin{tikzpicture}[>=Stealth,
    every node/.style={align=center, font=\small},
    root/.style={rectangle, rounded corners=1.5mm, draw=main, thick, fill=main!8},
    lv1/.style={rectangle, rounded corners=1.5mm, draw=second, thick, fill=second!8},
    lv2/.style={rectangle, rounded corners=1mm, draw=structurecolor, fill=structurecolor!8, text width=24mm}]
\coordinate (R) at (0,26mm);
\coordinate (A) at (-45mm,0);
\coordinate (B) at (-15mm,0);
\coordinate (C) at (15mm,0);
\coordinate (D) at (45mm,0);
\coordinate (A1) at (-45mm,-11mm);
\coordinate (A2) at (-45mm,-23mm);
\coordinate (B1) at (-15mm,-11mm);
\coordinate (B2) at (-15mm,-24mm);
\coordinate (B3) at (-15mm,-37mm);
\coordinate (C1) at (15mm,-11mm);
\coordinate (C2) at (15mm,-24mm);
\coordinate (C3) at (15mm,-37mm);
\coordinate (D1) at (45mm,-11mm);
\coordinate (D2) at (45mm,-24mm);
\coordinate (D3) at (45mm,-37mm);
\begin{scope}[on background layer]
\draw[thick] (R) -- (A); \draw[thick] (R) -- (B); \draw[thick] (R) -- (C); \draw[thick] (R) -- (D);
\draw (A) -- (A1); \draw (A) -- (A2);
\draw (B) -- (B1); \draw (B) -- (B2); \draw (B) -- (B3);
\draw (C) -- (C1); \draw (C) -- (C2); \draw (C) -- (C3);
\draw (D) -- (D1); \draw (D) -- (D2); \draw (D) -- (D3);
\end{scope}
\node[root] at (R) {数字经济指数 DEI};
\node[lv1] at (A) {宏观经济};
\node[lv1] at (B) {基础能力};
\node[lv1] at (C) {基础产业};
\node[lv1] at (D) {融合应用};
\node[lv2] at (A1) {三次产业\\增加值};
\node[lv2] at (A2) {信息消费\\规模};
\node[lv2] at (B1) {宽带与\\移动用户};
\node[lv2] at (B2) {云计算与\\大数据};
\node[lv2] at (B3) {物联网与\\终端};
\node[lv2] at (C1) {ICT 主营\\业务收入};
\node[lv2] at (C2) {电子信息\\进出口};
\node[lv2] at (C3) {互联网\\投融资};
\node[lv2] at (D1) {电子商务};
\node[lv2] at (D2) {互联网服务};
\node[lv2] at (D3) {"互联网+"\\融合};
\end{tikzpicture}
\caption{数字经济指数（DEI）指标体系示意（信通院，2017）}
\label{fig:dei}
\end{figure}

\subsection{两种测度方法的关系}

\important{增加值法与指数法并非替代关系，而是互补关系}：增加值法提供与 GDP 可比的绝对规模，用于国际比较与增长贡献核算；指数法提供多维度的相对评价，用于区域比较与政策评估。考试中的经典命题是"两种方法的适用场景比较"：当问题是"数字经济对 GDP 增长的贡献率是多少"时，须用增加值法；当问题是"两地数字经济发展水平孰高孰低"时，指数法更为适用。这一分工在实际报告中清晰可见：信通院既发布增加值规模，也发布 DEI 指数，前者用于说明总量，后者用于说明态势。

\begin{chapterbox}
\textbf{本章考点清单}

\begin{enumerate}
\item 数字经济定义的演进：Tapscott（1996）、OECD（2020）与 G20（2016）口径，三个"以"字结构的要素拆解 \important{【专硕·核心】}
\item 数字产业化（前四类）与产业数字化（第五类）的边界判定规则 \important{【专硕·核心】}
\item 中美数字经济占比的口径辨析（BEA 窄口径 vs 信通院宽口径）\important{【专硕·核心】}
\item "四化"框架的构成与相互关系 \important{【专硕·核心】}
\item 三要素（数据、算法、算力）的含义与相互关系 \important{【专硕·扩展】}
\item 三大定律的数学表达与推导：$N(t)=N_0 2^{t/T}$、$V\propto n^2$（及 $n\ln n$ 修正）、带宽增长及其与摩尔定律的速度比较 \important{【专硕·核心】}
\item 零边际成本与边际收益递增的经济含义（竞争定价失效与竞争均衡存在性条件的推导） \important{【专硕·核心】}
\item 概念谱系：信息经济、互联网经济、网络经济、平台经济、共享经济与数字经济的演进关系 \important{【专硕·核心】}
\item 发展三阶段与互联网泡沫案例 \important{【基础】}
\item 全球格局与中国"2522"政策框架 \important{【专硕·扩展】}
\item 增加值测度法四步、直接法与间接法、剥离比例 $s_d$ 的含义 \important{【专硕·核心】【学硕·热点】}
\item 指数测度法：极差标准化 $x_j'$ 与加权合成 $I=\sum w_j x_j'$ \important{【专硕·扩展】}
\item DEA-Malmquist 方法思想简介与中国数字经济指数 \important{【专硕·扩展】}
\item 测度三大难点：免费产品、质量调整、重复计算与剥离比例 \important{【专硕·扩展】}
\end{enumerate}
\end{chapterbox}

本章界定了数字经济的对象与特征，并给出了测度其规模的方法。第 2 章转入分析工具的准备，建立最优化与博弈论方法；第 3 章以网络外部性模型开启数字经济的微观机制分析，梅特卡夫定律的价值论将在那里被形式化。
